\chapter{Einleitung}\label{ch:intro}

\section{Motivation}\label{sec:motivation}

Suchalgorithmen und Indexstrukturen sind das Rückgrat moderner Hauptspeicher-Datenbanken;
trie-basierte Indizes und cache-line optimierte Verarbeitung stehen im Zentrum dieser Arbeit. Seit
dem Adaptive Radix Tree (ART) von Leis et al.~\cite{leis2013art} im Jahr 2013 hat sich die
Forschung in mehrere Richtungen ausdifferenziert: Trie-Höhen-Optimierung
(HOT~\cite{binna2018hot}), hybride B\textsuperscript{+}-Strukturen
(Masstree~\cite{mao2012masstree}, B\textsuperscript{2}-Baum~\cite{schmeisser2022b2tree}),
succinct-Repräsentationen (CoCo-Trie~\cite{boffa2024coco}, SuRF~\cite{zhang2018surf}) und
selbstoptimierende Ansätze (START~\cite{fent2020start}).

Parallel dazu hat sich ein eigener Forschungskorpus zu Cache-Bewusstheit, Prefetching und
Allokationsstrategien entwickelt. Heuristiken wie Cache-Sensitive Search Trees,
cache-oblivious Layouts und hierarchisches Prefetching adressieren jeweils eine separat betrachteter Aspekt
der Suchalgorithmus-Architektur. Aufgrund der fundamentalen Vorkommnisse auf cache-line optimierem
Verhalten basierender Code, ist eine Segmentierung der Suchalgorithmus-Hüllen-Bestandteile in Achsen sinnvoll.
Dies entspricht der rückwärts Abstraktion von bekannten Suchalgorithmus-Bestandteilen in wiederverwendbare
Module über die wichtigsten Paper des Standes der Technik hinweg.

Diese Arbeit schließt diese Lücke mit einer mehrschichtigen Cache-Engine-Architektur innerhalb
des sogenannten CacheEngine Frameworks
(ExecutionEngine $\to$ SearchEngine $\to$ SearchEngineType $\to$ SearchEngineType-Algorithmus-Permutation),
deren Bausteine als austauschbare
Entwurfsdimensionen eines großen Entwurfsraums (\emph{design space}) aufgefasst
werden~\cite{idreos2018periodic, idreos2018datacalculator}. Strukturell verschiedene
Suchverfahren --- von baumartigen Indizes bis zu nicht-baumartigen Hash-Tabellen, sowie 
Spezialeigenschaften wie cache-lines --- werden so
als Punkte desselben Entwurfsraums vergleichbar; nicht-baumartige Verfahren belegen dabei ein
eingeschränktes Subset (unterschiedlicher SearchEngineType) der Dimensionen mit denselben Suchalgorithmus-interfaces. 
Als Prüfling für die Stand-der-Technik-Zugehörigkeit dient PRT-ART (ein Probst-Redirect-Tree-/ART-Hybrid)
als eine abstrakte unvollständige Ansammlung neuartiger Suchalgorithmus-Bestandteile.

\section{Problemstellung}\label{sec:problem}

Bestehende cache-bewusste Index-Entwürfe optimieren ein oder wenige Aspekte isoliert und
berichten Ergebnisse, die sich über Veröffentlichungen hinweg schwer vergleichen lassen, da
jede Arbeit eigene Datenstrukturen, Workloads und Mess-Methodiken verwendet. Es fehlt ein
gemeinsamer Rahmen, der (i) solche Entwürfe in austauschbare Achsen zerlegt, (ii)
publizierte Algorithmen als explizite Konfigurationen rekonstruiert und (iii) sowohl
einzelne Achsen als auch ganze Algorithmen auf einer einheitlichen, fairen Basis vermessen kann.

\section{Forschungsfragen}\label{sec:rqs}
% Begriffsdefinition von SOTA und anderen Abkürzungen gehören in ein Glossar, sowie die Erklärung der Ränge in den Text
% Bezüglich der Forschungsfragen finden sich weitere Fragen in den Dokumenten der vorangegangenen Termine, wir müssen diese übernehmen
% In allen gesammelten Papern, die wir vollständig bezüglich ihrer Forschungsfragen ausnehmen müssen, finden wir hinweise, welche Aspekte im Fokus der cache-line Optimierung über verschiedene Suchalgorithmus-Bausteine hinweg, wir noch vergessen haben. Ich denke 3 Forschungsfragen sind zu wenig und wir sollten im Web nach Diplomarbeiten aus der Informatik suchen, um den Umfang besser abzuschätzen
\begin{description}
  \item[FF1 (Komposition)] Welche Algorithmus-Achsen eines cache-line-bewussten und damit cache-bewussten
        Trie-Index lassen sich kanonisch in einer Permutationsmatrix systematisieren, sodass
        Stand-der-Technik-Implementierungen als Profil-Konfigurationen rekonstruiert werden
        können? Lassen sich cache-line bewusste Messungen über multiple Aspekte von Suchalgorithmen mithilfe
        dieses Messsystems isolieren und leichter betrachten?
  \item[FF2 (Mess-Methodik)] Wie lässt sich pro Permutation ein fairer, reproduzierbarer
        YCSB-Workload-Lauf durchführen, der eine profilspezifische Workload-Auswahl (z.\,B.\
        ein Range-Scan-Profil $\to$ YCSB-E) erlaubt und dabei den Mess-Overhead über einen
        opt-in-Experiment-Modus minimiert? Welche Workloads gehören gerade für Suchalgorithmen zum Stand der Technik?
  \item[FF3 (Prüfling-Vergleich)] Wie schneiden die angebotenen Permutationen des PRT-ART als Prüfling und Testintegration             gegen die acht Rang-1-SOTA-Entwürfe (ART, HOT, Masstree, CoCo-Trie, START, B\textsuperscript{2}-Baum,
        Wormhole, SuRF) und die Rang-2/3-SOTA unter den drei Pflicht-Messreihen A (Prüfling
        vs.\ SOTA), B (SOTA-Vergleichsbasis) und C (Merge-Punkte) ab?
\end{description}

\section{Zielsetzung und Beiträge}\label{sec:contributions}

Die zentralen Beiträge dieser Arbeit sind:
% im Text vorhin habe ich die Schichten der Architektur nochmal nachgezogen, weiterhin ist der Code der Diplomarbeit der formale Bediener der cache-engine, welche sich den PRT-ART als Prüfling herunterläd und diesen per Metaprogrammierung gegen die Kernbestandteile der cache-engine kompiliert und als ABI-stabile Lebewesen ausliefert und über das Prüfdock zur Messung jagt. Die Zielsetzung ist korrekt, aber noch etwas unscharf bzw. da fehlen etliche Punkte und wir können gerne die Zielstellung scharf zeichnen.
\begin{itemize}
  \item Eine dreischichtige Architektur (CacheEngine $\to$ ExecutionEngine $\to$
        SearchEngine) mit ABI-stabilen C++23-Modul-Schnittstellen und variadisch
        templatierten Hybrid-API-Spezialisierungen (1~Parameter $\Rightarrow$
        \texttt{std::vector}-API; 2~Parameter $\Rightarrow$ \texttt{std::map}-API; $N>2
        \Rightarrow$ \texttt{std::map<K,\,std::tuple<V$_1$\ldots V$_N$>>}).
  \item Eine neunzehn-Achsen-Permutationsmatrix mit \texttt{std::variant}-Bausteinen und
        Compile-Time-Fallback zwischen Prüfling-Stack und Stand-der-Technik-Stack.
  \item Dreißig SOTA-Algorithmus-Profile als persistierte XML-Konfigurationen (8~Rang-1 +
        22~Rang-2/3), die der \texttt{ExperimentDriver} automatisch in vorkompilierte
        C++23-Module überführt.
  \item PRT-ART als Prüfling mit acht eigenen Bausteinschichten (4+2-Pool-Allokator,
        optimistisches Lock-Coupling mit reservierten Blöcken, MultiLevel-Layout,
        pfadorientiertes Prefetch, DensityTracker, Hypothesen-Metriken H1/H2/H3,
        Inline/External/ChainRef-ValueHandles, Signaling-Bit-Serialisierung).
  \item Ein Mess-Treiber mit Defined-/Full-Modus-Unterscheidung, profilbewusstem
        Workload-Routing und drei Mess-Granularitäten (einzelne Achsen, Interface-Operationen,
        ganze Algorithmen).
  \item Die Auslieferung des besten Gesamtalgorithmus --- ausgewählt aus den feingliedrigen
        Messergebnissen --- als eigenständige, ABI-stabile Binary hinter dem einheitlichen
        Suchalgorithmus-Interface.
\end{itemize}

\section{Aufbau der Arbeit}\label{sec:structure}
% Es sind nicht nur Such- und Allokator-Cluster. Es gibt mehrere Achsen Repositories, die thematisch mehrere Achsen abbilden und die Sub-Achsen sollten wir auch nicht vergessen, die derzeit dokumentiert sind. Da die Diplomarbeit durchaus eine Arbeit mit 60 bis 80 Seiten eigenem Volltext sein darf, ist genug Platz für Details der mehrstufigen Untergliederung der Architektur. Derzeit fehlt aber hauptsächlich im Stand der Technik ein eigenes Unterkapitel, was sich mit dem nicht-technischen Bereich des Suchalgorithmus Architektur-Designs beschäftigt, weil unser Framework eine Architektur-Design-Erfindung mithilfe von fortgeschrittenen Design-Pattern zur Lösung eines technischen Problems darstellt, welches bisher immer mit "bloßen Händen" gelöst wurde, wir konstruieren aber vor der Arbeit erst eine Schaufel -> lösen das Problem also indierekt auf einem nicht sofort ersichtlichen Weg. Der Stand der Technik ist komplett richtig, muss aber durch Konteption und Voraussetzungen von Suchalgorithmen beginnen, was eine Webrecherche erfordert und natürlich professionelle wissenschaftliche Zitate. Wir bewegen uns also parallel zur Forschungsfeld der Datenbanken im Bereich des Software-Engineerings an der TU Dresden und deren Zielen (Webrecherche bitte).
Kapitel~\ref{ch:fundamentals} führt in die nötigen Grundlagen ein.
Kapitel~\ref{ch:sota} fasst den Stand der Technik über sechs Such-Cluster (A--F) und einen
Allokator-Cluster zusammen. Kapitel~\ref{ch:concept} stellt das Achsen-Bibliotheks-Framework,
die geschichtete Architektur, die ABI-Schnittstelle und die neunzehn-Achsen-Permutationsmatrix
vor. Kapitel~\ref{ch:impl} dokumentiert die Implementierung über die drei Repositorien.
Kapitel~\ref{ch:methodology} beschreibt die Mess-Methodik mit den drei Pflicht-Messreihen, den
drei Mess-Granularitäten, den realen Datensätzen und den Fairness-Regeln. Kapitel~\ref{ch:results} präsentiert und diskutiert die
Ergebnisse einschließlich der Hypothesen H1--H3. Kapitel~\ref{ch:conclusion} schließt mit
Zusammenfassung und Ausblick.
